\paragraph{不同师生对上的推理迁移。}
表~\ref{tab:historyendpoints}展示两组Step200 Avg@8差值为正的历史案例。
Qwen3-1.7B-Base学生与Rethinking OPD\citep{rethinking}发布的公开GRPO教师配对时，Block3的四项Avg@8均高于Token：
MATH500从45.85升至54.75，AIME24从5.42升至8.75，AIME25从2.50升至5.42，
AMC23从21.84升至27.26。这些观察推动我们进一步研究超出逐token更新的局部反馈共享。
DeepSeek/JustRL 1.5B对照的终点差值也为正，MATH500和AMC23的变化小于AIME。

这两组是正向案例展示，不是全部跨模型对照；其他历史师生对存在更有利于Token的结果。
附录\ref{app:history}表~\ref{tab:historyendpointsfull}保留完整四组终点对照及协议背景，
随后给出全部checkpoint结果。

\begin{table}[htbp]
\centering\small
\caption{选取的历史Step200 Avg@8正向案例（\%）。仅在每个学生分层内比较方法；包括负结果的完整对照见表~\ref{tab:historyendpointsfull}。}
\label{tab:historyendpoints}
\begin{tabular}{@{}llrrrr@{}}
\toprule
Student & Method & MATH500 & AIME24 & AIME25 & AMC23 \\
\midrule
Qwen 1.7B & Token & 45.85 & 5.42 & 2.50 & 21.84 \\
Qwen 1.7B & Block3 & 54.75 & 8.75 & 5.42 & 27.26 \\
\addlinespace
DeepSeek 1.5B & Token & 84.70 & 39.58 & 29.17 & 73.34 \\
DeepSeek 1.5B & Block3 & 84.75 & 43.75 & 30.83 & 73.80 \\
\addlinespace
\bottomrule\end{tabular}

\end{table}

\paragraph{官方Qwen3-8B到1.7B指令版。}
表~\ref{tab:qwen17endpoint}在原生非thinking chat、训推输入一致的设置下比较两种方法。
相对Token，Block3的Step200 Avg@8在MATH500、AIME24、AIME25和AMC23上分别高
$+0.10$、$+2.50$、$+0.83$和$+0.30$个百分点。
两种方法的四项Avg@8均超过原始学生。
方法差值的逐项95\%配对题目区间均包含零，因此正向点估计尚不能在该端点上明确区分统计优势。

Pass@8区分平均成功率与题目覆盖。
Block3相对Token的差值分别为$+0.20$、$-6.67$、$0.00$和$-1.20$个百分点，
且两种方法的MATH500覆盖率均低于原始学生的91.20\%。
因此，平均准确率提高可以与八次尝试覆盖率下降同时出现。
附录\ref{app:instruct}报告原始学生、全部八个训练权重和对应行为指标。
\begin{table}[t]
\centering\small
\caption{官方Qwen3-8B到Qwen3-1.7B指令版：Step200分数（\%）及Block3减Token差值（百分点）。}
\label{tab:qwen17endpoint}
\begin{tabular}{lrrrr}
\toprule
Benchmark & Original & Token & Block3 & $\Delta$ [95\% CI] \\
\midrule
\multicolumn{5}{c}{Avg@8} \\
MATH500 & 71.45 & 72.12 & 72.22 & +0.10 [-1.18, 1.38] \\
AIME24 & 10.00 & 14.17 & 16.67 & +2.50 [-2.50, 8.33] \\
AIME25 & 10.42 & 12.08 & 12.92 & +0.83 [-3.75, 5.83] \\
AMC23 & 40.51 & 43.07 & 43.37 & +0.30 [-3.16, 3.92] \\
\midrule
\multicolumn{5}{c}{Pass@8} \\
MATH500 & 91.20 & 89.60 & 89.80 & +0.20 [-2.00, 2.40] \\
AIME24 & 26.67 & 43.33 & 36.67 & -6.67 [-20.00, 6.67] \\
AIME25 & 23.33 & 30.00 & 30.00 & +0.00 [-16.67, 16.67] \\
AMC23 & 69.88 & 75.90 & 74.70 & -1.20 [-9.64, 7.23] \\
\bottomrule
\end{tabular}
\par\smallskip\footnotesize Original为本实验训练前的指令版学生，不是Base模型。方括号为10,000次配对题目重采样的逐点95\%区间；仅一个训练seed，未校正checkpoint或任务选择。
\end{table}


\paragraph{Qwen3-4B Base与公开GRPO教师。}
采用匹配completion提示的对照呈现不同的终点模式（表~\ref{tab:qwen4endpoint}）。
Step200时，Block3的Avg@8差值依次为$-0.52$、$-2.92$、$+0.42$和$-3.01$个百分点。
Token在三项任务上更高，但逐项95\%配对题目区间均跨零，既未确立优越性，也未证明等效。
AIME24中，Block3的Avg@8更低，却有多一道题至少答对一次，Pass@8为33.33\%，Token为30.00\%。
综合这些成对结果，信用粒度会影响学习轨迹，其收益与师生设置及所考察的指标有关。
% Generated from accepted results; requires booktabs.
\begin{table*}[t]
\centering
\scriptsize
\setlength{\tabcolsep}{4pt}
\caption{Qwen3-4B Step200 配对评测 (seed21; 队列 COMPLETE)。}
\label{tab:qwen4endpoint}
\begin{tabular}{lrrrrrr}
\toprule
 & \multicolumn{3}{c}{Avg@8} & \multicolumn{3}{c}{Pass@8} \\
\cmidrule(lr){2-4}\cmidrule(lr){5-7}
基准 & Block3 & Token & $\Delta$ [95\% CI] & Block3 & Token & $\Delta$ [95\% CI] \\
\midrule
MATH500 & 75.40 & 75.92 & -0.52 [-1.57, 0.55] & 90.60 & 91.00 & -0.40 [-2.40, 1.60] \\
AIME24 & 13.33 & 16.25 & -2.92 [-7.08, 0.83] & 33.33 & 30.00 & +3.33 [-6.67, 13.33] \\
AIME25 & 13.75 & 13.33 & +0.42 [-2.92, 3.75] & 30.00 & 30.00 & +0.00 [-13.33, 13.33] \\
AMC23 & 46.23 & 49.25 & -3.01 [-6.17, 0.30] & 74.70 & 77.11 & -2.41 [-10.84, 6.02] \\
\bottomrule
\end{tabular}
\par\smallskip\footnotesize 分数为百分比；$\Delta$为Block3减Token的百分点差。方括号为各基准内成对问题bootstrap的逐点95\%区间，非多训练种子CI；NA表示缺失。
\end{table*}


\IfFileExists{\qwenfigures/qwen4_accuracy.pdf}{%
\begin{figure}[t]
\centering
\includegraphics[width=\linewidth]{\qwenfigures/qwen4_accuracy.pdf}
\caption{当前匹配4B对照的四个保存权重。各benchmark独立报告Avg@8，使用独立纵轴范围，同一子图内两种方法共享范围；连线连接实测权重，不代表额外评测。每种方法仅一个训练seed。}
\label{fig:qwen4accuracy}
\end{figure}
}{}
